\SourceLeaf{459}
\noindent Dividiert man diese beiden Gleichungen durcheinander und zieht dann die Quadratwurzel, so ergibt sich
\[
\frac{\sin\tfrac12 a^3}{\cos\tfrac12 a}.\frac{\cos\tfrac12 b}{\sin\tfrac12 b^3}
=\frac{\sin(A+\omega)\sin(A-\tfrac12\omega)^2}{\sin(B+\omega)\sin(B-\tfrac12\omega)^2}.
\]
\noindent Man kann diese Gleichung auch in die Form bringen
\[
\frac ab=\frac{\sin A}{\sin B}.\mkern6mu{}^3\mkern-4mu\BareRoot{D},
\]
\Needspace{10\baselineskip}
\noindent wobei zur Abkürzung $D$ anstelle von
\[
\begin{aligned}
&\frac{a^3\cos\tfrac12 a}{8\sin\tfrac12 a^3}.\frac{8\sin\tfrac12 b^3}{b^3\cos\tfrac12 b}\, .\\[5pt]
&\quad\frac{\sin(A+\omega)\sin(A-\tfrac12\omega)^2}{\sin A^3}.\frac{\sin B^3}{\sin(B+\omega)\sin(B-\tfrac12\omega)^2}
\end{aligned}
\]
\noindent geschrieben ist.

Diese Formel ist streng richtig. Man sieht aber sofort, dass, wenn $a,b,c$ sehr klein sind und als Größen erster Ordnung betrachtet werden, jeder der vier Faktoren, aus denen $D$ zusammengesetzt ist, von der Einheit nur um Größen vierter Ordnung abweicht.

Nach allgemeineren Prinzipien ist dieser Gegenstand in meinen \textit{Disquisitiones generales circa superficies curvas.} behandelt und auf die Dreiecke ausgedehnt, die auf beliebigen krummen Flächen zwischen kürzesten Linien gebildet werden.

